\documentclass[11pt]{article}
\usepackage[review]{acl}
\usepackage{times}
\usepackage{latexsym}
\usepackage[T1]{fontenc}
\usepackage[utf8]{inputenc}
\usepackage{microtype}
\usepackage{inconsolata}
\usepackage{graphicx}

\title{RouteBench: Evaluating Least-Authority Capability Routing in Code Agents}

\author{Anonymous ACL Submission}

\begin{document}
\maketitle

\begin{abstract}
Code agents increasingly choose among multiple tool routes that can complete the same
task but expose different authority. Existing evaluations often collapse task success
and safety into a single outcome, making it difficult to distinguish inability from
unsafe route selection. We present RouteBench, a deterministic benchmark for evaluating
whether an agent can select a functionally successful route with globally least
authority from a bounded capability graph. RouteBench separates functional completion,
policy compliance, and least-authority route selection, and verifies outcomes through a
trusted stateful simulator rather than textual traces. Its formal contract uses manually
authored schema-v4 scenarios, hidden route certificates, typed tools, and held-out
template families. We describe the benchmark design, threat model, evaluation protocol,
and the data boundary required for a future evaluation-calibrated reinforcement-learning
defense. Results and defense experiments are intentionally reported only after the
corresponding frozen runs are complete.
\end{abstract}

\section{Introduction}

Language-model agents increasingly interact with databases, files, identity systems,
cloud resources, and external services through typed tools. A task can therefore have
several functionally valid routes with different authority footprints. An agent may
complete the task while selecting an unnecessarily broad capability, or it may follow a
functionally plausible route that realizes a forbidden effect. These cases are not
captured by task success alone.

We introduce RouteBench, a benchmark for this capability-routing problem. Given a
bounded directed graph of typed actions, an agent must select a route that completes the
ordinary task while minimizing authority under a frozen evaluator contract. The
benchmark distinguishes three outcomes: whether the task completes, whether the route
is policy compliant, and whether it matches the unique least-authority optimum. Unsafe
branches are verified by deterministic state assertions and are reported separately
from planner capability failures.

RouteBench is designed around four requirements. First, the evaluator must verify
effects in a clean state rather than trust a textual plan. Second, the public task and
graph must not reveal the evaluator-owned optimum. Third, benchmark cases must cover
distinct permission dimensions, tool surfaces, route topologies, and attack mechanisms.
Finally, formal results must remain separated from any later policy training or defense
calibration.

Our contributions are:
\begin{itemize}
  \item a formal capability-routing task with typed tools, bounded candidate graphs, and
        hidden least-authority certificates;
  \item a deterministic evaluation protocol that separates functional success, policy
        compliance, optimal-path success, and attack-path selection; and
  \item a data and split contract that supports future defenses, including constrained
        reinforcement learning, without leaking held-out test cases into training.
\end{itemize}

\section{Problem Definition}

Let a scenario provide a task $T$, an initial state $s_0$, and a bounded directed graph
$G=(V,E)$. Each edge represents one typed tool action and declares an authority atom.
A route $p$ is executable when its dependencies hold at every step. The trusted
simulator maps a route to a terminal state and an authority footprint.

The evaluator accepts a route only when it satisfies the functional verifier and policy
verifier. Among accepted routes, the primary objective is the evaluator-defined
least-authority route, with action count used as a secondary objective where required.
The model receives the public task projection, tools, graph, and objective description;
it does not receive the reference route, hidden certificate, fixture, or outcome labels.

\section{Benchmark Design}

\subsection{Scenario Construction}

Formal cases are manually authored and use schema v4. Each case contains a public task
and graph, a private fixture, typed tool definitions, functional and policy verifiers,
and evaluator-owned safe and attack certificates. Construction checks establish that a
safe functional route exists, that the least-authority certificate is unique, and that
attack certificates produce the declared forbidden effects.

The taxonomy covers eight forbidden-effect families, eight permission dimensions, ten
typed tool surfaces, five route topologies, five route-length bands, and six attack
mechanism labels including fixed control, parameter substitution, equivalent-tool
substitution, decomposition, delayed reordering, and authorization-state
manipulation.

\subsection{Evaluation Outcomes}

For each attempt, the simulator records functional success, policy compliance,
attack success, certificate matching, public Pareto optimality, and public exact
optimality. The primary formal endpoint is \texttt{OptimalPathSuccess}. A case is
considered capable only when at least three of five declared seeds achieve that endpoint;
backend failures remain infrastructure failures and are not converted into model
failures.

\section{Evaluation Protocol}

The formal inventory contains 776 held-out primary cases distributed across frozen
template families. Each case is evaluated with five seeds at temperature zero using the
same typed interface and trusted simulator. Run manifests record the Git revision,
dependency and planner hashes, model identity, weight hash, scenario digests, and
generation parameters. Raw prompts, responses, fixtures, graphs, and hidden
certificates are excluded from aggregate artifacts.

\section{Toward Evaluation-Calibrated Defense}

RouteBench is also intended to support a separate defense study. Development traces
can be converted into state-action-outcome records for a policy that mediates typed
capability requests. A future constrained offline-RL policy may choose to allow,
constrain, escalate, or deny an action using the task, current capability state, route
history, and data-flow labels. Formal test cases must remain excluded from training,
reward design, threshold selection, and policy revision. This separation makes the
defense result a held-out generalization test rather than a benchmark-specific rule.

\section{Discussion and Limitations}

RouteBench uses deterministic simulated tools and therefore does not by itself establish
security against arbitrary real-world APIs or untyped shell behavior. The benchmark
measures route selection and trusted state effects; it does not claim that a model's
internal reasoning is observable or safe. These limits motivate reporting both the
capability endpoint and the authorization outcomes, and motivate external validation of
future defenses on independently authored workflows.

\section{Conclusion}

RouteBench frames agent safety as least-authority capability routing. Its separation of
task completion, policy compliance, and route optimality provides a basis for measuring
where agents fail and for evaluating defenses without conflating capability failure
with security success. The current manuscript establishes the benchmark contract; all
empirical claims will be tied to frozen, provenance-recorded runs.

\bibliography{custom}
\bibliographystyle{acl_natbib}

\end{document}
